\begin{answer}
\subsection*{i)}
 
\[
    \text{Let } X = \begin{bmatrix} \dots (x^{(1)})^T \dots \\ \dots (x^{(2)})^T \dots \\ \vdots \\ \dots (x^{(n)})^T \dots \end{bmatrix}, \quad Y = \begin{bmatrix} y^{(1)} \\ y^{(2)} \\ \vdots \\ y^{(n)} \end{bmatrix}
\]

\vspace{0.5em}
As our weights are just multiplied in the summation of J($\theta$), it must be a diagonal matrix, defined as:

\[
    W = \begin{bmatrix} w_1 & & 0 \\ & \ddots & \\ 0 & & w_n \end{bmatrix}
    \]
    
Let the error E be:
    \[
    E = X\theta - Y = \begin{bmatrix} (x^{(1)})^T \theta - y^{(1)} \\ \vdots \\ (x^{(n)})^T \theta - y^{(n)} \end{bmatrix}
    \]
    
The cost function $J(\theta)$ is:

    \begin{align*}
        (1) \quad J(\theta) &= \sum_{i=1}^{n} w^{(i)} ((x^{(i)})^T\theta - y^{(i)})^2 = \sum_{i=1}^{n} w^{(i)} e_i^T e_i \\
        &= \frac{1}{2} E^T W E = \frac{1}{2} (X\theta - Y)^T W (X\theta - Y)
    \end{align*}
    \begin{center}
        \small{(The scaling factor of 0.5 can be omitted, as it doesn't change the landscape of the loss. So it does fit the equation in the problem statement.)}
    \end{center}
    
    % Section II
    \subsection*{ii)}
    \textbf{Derivative of weighted $J(\theta)$:}
    \begin{align*}
        \nabla_\theta J(\theta) &=\nabla_\theta \frac{1}{2} ((X\theta - Y)^T W (X\theta - Y)) \\
        &= \nabla_\theta \frac{1}{2}(\theta^T X^T W X \theta - \theta^T X^T W Y - Y^T W X \theta) \\
        &= \nabla_\theta \frac{1}{2}(\theta^T X^T W X \theta) - 2 \nabla_\theta (\theta^T X^T W Y) \\
    \end{align*}
    With $\nabla_z z^T A z = 2Az$ for $z=\theta$ and $A=X^T W X$, we can simplify even further:
\begin{align*}
        \nabla_\theta J(\theta) &= \frac{1}{2} (2X^T W X \theta - 2X^T W Y) \\
        &= X^T W X \theta - X^T W Y
        \end{align*}
        
    \textbf{Find minimum of $J(\theta)$:}
    \begin{align*}
        \nabla J(\theta) &= 0 \\
        0 &= X^T W X \theta - X^T W Y \\
        X^T W X \theta &= X^T W Y \\
        % Solving for \theta gives the final expression for weighted least squares:
        \theta &= (X^T W X)^{-1} X^T W Y = \arg\min_\theta J(\theta) \\
    \end{align*}

    \subsection*{iii)}
    The Maximum Likelihood Estimate requires us to find, given the data we observed, 
the set of parameters \(\theta\) that makes it the most likely.
For this, we assume that our dataset underlying distribution follows a normal distribution
with mean \(\theta^T x^{(i)}\) and variance \(\sigma^{(i)}\).
For this, we can use the so-called likelihood function, giving us the likelihood
of observing a specific set of data for that parameter \(\theta\). \\

As each sample is independent we can also say:
\begin{align*}
    L(\theta) &= p(y|X; \theta) = \prod_{i=1}^{k} p(y^{(i)}|x^{(i)}; \theta) = \prod_{i=1}^{k} \frac{1}{\sqrt{2\pi}\sigma^{(i)}} \exp\left(-\frac{(y^{(i)} - \theta^T x^{(i)})^2}{2(\sigma^{(i)})^2}\right)
    \end{align*}
    
    To make finding argmax $L(\theta)$ easier, we can use the log-probability $l(\theta)$:
    \begin{align*}
    l(\theta) &= \log(L(\theta)) = \sum_{i=1}^{n} \log(p(y^{(i)}|x^{(i)}; \theta)) \\
    &= \sum_{i=1}^{n} \left( \log\left(\frac{1}{\sqrt{2\pi}\sigma^{(i)}}\right) - \frac{(y^{(i)} - \theta^T x^{(i)})^2}{2(\sigma^{(i)})^2} \right)
    \end{align*}

    To maximize, we need to minimize the second (and only) term that depends on \(\theta\):
\begin{align*}
    & \underset{\theta}{\text{argmin}} \left( \sum_{i=1}^{n} \frac{1}{2(\sigma^{(i)})^2} (y^{(i)} - \theta^T x^{(i)})^2 \right) \\
    = & \underset{\theta}{\text{argmin}} \left(\frac{1}{2} \sum_{i=1}^{n} \frac{1}{(\sigma^{(i)})^2} (y^{(i)} - \theta^T x^{(i)})^2 \right)
    \end{align*}

    The above expression looks exactly like the weighted least squares loss from (i).
    The \(w^{(i)}\) term is equivalent to \(\frac{1}{(\sigma^{(i)})^2}\) in the MLE.

\end{answer}
